\section{Methods --- TOS-CMI as measurement $\to$ certified control}

Four layers (localize $\to$ protect $\to$ gate $\to$ refuse) plus
the evaluation protocol. The method is a \textbf{measurement-and-certification} procedure, not an
always-on regularizer.

\subsection{Problem setup and conditional leakage}
We observe a representation $Z = f(X)$ of EEG trials with task labels $Y$ and a domain (subject) variable $D$.
Conditional domain leakage is the residual dependence of $Z$ on $D$ given the task, $I(Z;D|Y)$: a representation
can be domain-invariant in aggregate yet still encode subject identity \emph{within} each class. Unlike a global
marginal penalty, we treat $I(Z;D|Y)$ as something to first \textbf{localize} in $Z$ and then \textbf{delete only if
certified safe} for the task. Throughout, the held-out target domain is never used for selection, gating,
or calibration --- only for final reporting (\S2.5).

\subsection{Score-Fisher localization}
We localize the directions of $Z$ that are domain-rich but task-light using a \textbf{conditional score-Fisher}
construction. Let $G_Y$ be the label score-Fisher matrix (outer products of $\nabla_z \log p(Y|z)$) and $G_{D|Y}$ the
conditional-domain score-Fisher matrix ($\nabla_z \log q(D|z,Y)$, formed on the class-conditional measure). Both are
estimated by cross-fitting (out-of-fold scores) under a whitening metric $M = (\Sigma_W + \epsilon \Sigma_{ref})^{-1}$. The
candidate nuisance subspace $V_D$ is read off the leading generalized eigenvectors of

\begin{equation}
G_{D|Y} v = \rho (G_Y + \eta M) v,
\end{equation}

i.e. directions that maximize conditional-domain score energy relative to task score energy (large $\rho$).
This is a \textbf{score-Fisher proxy} for conditional leakage --- second-order in the log-likelihood scores, so it
detects covariance/synergy structure that a first-moment (mean-scatter) statistic is blind to --- not the
exact conditional mutual information [C1]. The number of retained directions nDcand is chosen by an
eigengap/$\rho$-threshold rule.

\begin{figure}[t]
\centering
\fbox{\parbox{0.92\linewidth}{\centering Pipeline schematic placeholder: score-Fisher localization $\rightarrow$ direct-sum projection $\rightarrow$ task-risk / domain-gain gates $\rightarrow$ delete or abstain.}}
\caption{Overview of the TOS-CMI measurement-to-control framework: localize a domain-rich, task-light subspace; delete it with a direct-sum projector; gate on a conditional task-risk certificate; and \emph{delete or refuse}.}
\label{fig:pipeline}
\end{figure}

\subsection{Direct-sum projection and the low-rank deletion diagnostic}
Given $V_D$ we delete it with an \textbf{M-oblique direct-sum projector} $P_N$ onto $V_D$ along the task carrier, and
study the complement $RZ = (I - P_N) Z$. The projector is constructed to satisfy the direct-sum conditions

\begin{equation}
R V_D = 0 \quad (\text{the nuisance subspace is removed}) \quad \text{and} \quad R T = T \quad (\text{the task carrier is preserved}),
\end{equation}

rather than mere orthogonality. We then \textbf{measure removability} by decoding $D$ and $Y$ from $RZ$ (linear and
MLP probes) and comparing against (i) the full $Z$ and (ii) a same-k \emph{random} removal control. The deletion
is informative only if it reduces domain decode well below random-k while leaving task decode unchanged.
Crucially, \textbf{direct-sum geometry is necessary for algebraic task preservation but not sufficient for
conditional task safety} [C2]: a domain-rich direction can still carry conditional task information, which
the geometry cannot see --- motivating the gate.

\subsection{Conditional task-risk gate and certified refusal}
Before deletion is permitted, we certify that removing $V_D$ does not destroy task-relevant information by
bounding the conditional task risk it carries,

\begin{equation}
\Delta_Y(k) = I(Y; P_k Z | (I - P_k) Z),
\end{equation}

the task information in the deleted subspace beyond what the complement already holds. $\Delta_Y$ is estimated by a
\textbf{cross-fitted, one-step plug-in log-ratio} estimator with a power certificate (a minimum-detectable-effect
floor); a weak nested critic is insufficient because it can \textbf{unsafe-accept} deletions whose true $\Delta_Y$ is
large (Fig 2B). The gate combines an upper confidence bound on $\Delta_Y$ (task safety) with a lower bound on the
domain-gain (that the subspace is genuinely domain-rich). \textbf{Certification is allowed to abstain, and
abstention is a valid control decision}: if the gate cannot certify that deletion is both safe ($\Delta_Y$ small)
and useful (domain-rich) at the available sample size, the correct action is the \textbf{identity map}, not a
weaker invariance penalty [C3]. On synthetic controls this certified gate produces zero unsafe-accepts but
is deliberately conservative (Fig 2D) --- an honest negative for default-on deletion, by design.

\subsection{Frozen-feature EEG evaluation protocol}
We evaluate on BCI-IV-2a (BNCI2014\_001) under leave-one-subject-out (LOSO), domain = subject. For each
fold we train a backbone, \textbf{freeze} it, and dump its penultimate latent $Z$ on the source subjects; all
diagnostics (\S2.2--2.4) run on the frozen $Z$, and the held-out subject is used only for reporting target/LOSO
accuracy. We compare two representations to test generality: TSMNet (LogEig/SPD tangent latent, $d_z=210$)
and EEGNet (convolutional penultimate latent, $d_z=16$). As the standard baseline for ``global'' conditional
invariance we sweep an LPC penalty (CE + $\lambda \cdot I(Z;D|Y)$) over $\lambda$, with per-epoch instrumentation (task CE,
feature norm, penalty, encoder gradient, effective rank) to diagnose \emph{how} it changes the representation
(\S4.1). The instrumentation is read-only and flag-gated (default off), with verified zero effect on
training dynamics. Probes (subject/task decode) use a 50/50 trial split with linear and MLP heads; the
random-k removal control deletes the same number of M-orthonormal directions as $V_D$.

Note (frozen-feature scope): this protocol measures whether \emph{removing} localized leakage from a fixed
representation helps; it does not retrain the encoder with a certified penalty. The end-to-end variant is
out of scope here and flagged in \S6.3.
